% Source pp. 135-143 (Lectures13-18.pdf).
\section[Renormalization of Composite Fields (\texorpdfstring{$\varphi^4$}{phi**4} theory)]{Renormalization of Composite Fields (\texorpdfstring{$\varphi^4$}{phi**4} theory)}

% insertion vertices
\newcommand{\insv}[2][]{\dia[0.4]{\draw[prop,wavy] (-2.2,0) -- (0,0);
  \draw[prop] (0,0) -- (1.4,1.4); \draw[prop] (0,0) -- (1.4,-1.4); #2 \node[font=\tiny] at (-2.6,0.45) {#1};}}

In many cases one may be interested in correlation functions involving composite
fields, such as powers of $\varphi^n(x)$, or energy-momentum tensor. Such
correlation functions can be computed perturbatively, as the power series in the
coupling constant $\lambda$. However new renormalizations, in addition to the
basic renormalizations of the parameters $m^2$, $\lambda$ and the field
$\varphi$, have to be made to make such correlation functions finite. Let us
demonstrate this for the simplest case of the composite field
\[
\varphi^2(x)
\]
in the $\varphi^4$ theory.

Consider the 3-point connected correlation function
\[
\Bigl\langle\frac12\,\varphi^2(x)\,\varphi(x_1)\,\varphi(x_2)\Bigr\rangle_c
\]
where $\varphi(x)$ stands for the renormalized field, $\varphi=Z^{-1/2}\varphi_0$,
$Z=1+\delta Z$, with $Z$ chosen according to some scheme, say SII. Diagrams can
be used to describe perturbative contributions to this correlation function if
one introduces additional ``external'' vertex
\[
\insv[$\varphi^2$]{\node[vtx,minimum size=4.5pt] at (0,0) {};}=1
\]
associated with the insertion $\frac12\varphi^2$. As before, it is useful to
switch to the momentum representation, and introduce the associated amputated
vertex
\begin{equation}\label{eq:Gphi2}
-\Gamma^{(\varphi^2,2)}(p_1,p_2)=\frac{W^{(\varphi^2,2)}(p_1,p_2)}{W(p_1)\,W(p_2)}\,,
\end{equation}
where
\[
(2\pi)^4\delta^{(4)}(q-p_1-p_2)\,W^{(\varphi^2,2)}(p_1,p_2)=\text{f.t.}\,\Bigl\langle\frac12\,\varphi^2(x)\,\varphi(x_1)\varphi(x_2)\Bigr\rangle_c\,,
\]
with $q,p_1,p_2$ being the momenta corresponding to $x,x_1,x_2$, respectively.
It is given by the sum of all one-particle irreducible diagrams
\begin{multline*}
-\Gamma^{(\varphi^2,2)}(p_1,p_2)=
\underset{(a)}{\diac[0.4]{\draw[prop,wavy] (-2,0) -- (0,0); \draw[prop,densely dashed] (0,0) -- (1.2,1.2); \draw[prop,densely dashed] (0,0) -- (1.2,-1.2); \node[vtx] at (0,0) {};}}
+\underset{(b)}{\diac[0.4]{\draw[prop,wavy] (-2.8,0) -- (-1.6,0); \draw[prop] (-0.8,0) ellipse (0.8 and 0.5);
  \draw[prop,densely dashed] (0,0) -- (1.2,1.2); \draw[prop,densely dashed] (0,0) -- (1.2,-1.2); \node[vtx] at (-1.6,0) {}; \node[vtx] at (0,0) {};}}
+\underset{(c)}{\diac[0.4]{\draw[prop,wavy] (-4.4,0) -- (-3.2,0); \draw[prop] (-2.4,0) ellipse (0.8 and 0.5); \draw[prop] (-0.8,0) ellipse (0.8 and 0.5);
  \draw[prop,densely dashed] (0,0) -- (1.2,1.2); \draw[prop,densely dashed] (0,0) -- (1.2,-1.2); \foreach \x in {-3.2,-1.6,0} \node[vtx] at (\x,0) {};}}\\
+\underset{(d)}{\diac[0.4]{\draw[prop,wavy] (-2.8,0) -- (-1.6,0); \coordinate (T) at (-0.25,0.6); \coordinate (B) at (-0.25,-0.6);
  \draw[prop] (-0.8,0) circle (0.8); \draw[prop] (T) .. controls (-0.55,0) .. (B);
  \draw[prop,densely dashed] (T) -- ++(45:1.2); \draw[prop,densely dashed] (B) -- ++(-45:1.2); \foreach \p in {(-1.6,0),(T),(B)} \node[vtx] at \p {};}}
+\underset{(e)}{\diac[0.4]{\draw[prop,wavy] (-2.8,0) -- (-1.6,0); \draw[prop] (-0.8,0) ellipse (0.8 and 0.5);
  \draw[prop,densely dashed] (0,0) -- (1.2,1.2); \draw[prop,densely dashed] (0,0) -- (1.2,-1.2); \node[vtx] at (-1.6,0) {}; \node[otimesv,minimum size=5pt] at (0,0) {};}}
+\cdots
\end{multline*}
Most of these diagrams are singular at $d=4$, and a quick look reveals that these
singularities can not be eliminated by renormalizations of the parameters $m^2$,
$\lambda$, and the field $\varphi$. This is already seen in the order
$\sim\lambda$. All the counterterm vertices that we have introduced already show
up at the order $\lambda^2$ (diagram (e)) and higher. But the diagram (b), which
is $\sim\lambda$, is singular. Indeed its contribution is
\[
-\Gamma^{(\varphi^2,2)}(p_1,p_2)=1-\frac\lambda2\,I(p_{12}^2)+O(\lambda^2)
\]
where $I(p^2)$ is the same integral \eqref{eq:Ip2} which we encountered
previously. We know that
\[
I(p_{12}^2)=I(0)+I_r(p_{12}^2)\,,
\]
where $I_r(p^2)$ has finite limit at $d=4$, while $I(0)$ diverges as $1/\epsilon$.
In fact this shows what we need to do. We write
\[
-\Gamma^{(\varphi^2,2)}(p_1,p_2)=\Bigl(1-\frac\lambda2\,I(0)+O(\lambda^2)\Bigr)\Bigl(1-\frac\lambda2\,I_r(p_{12}^2)+O(\lambda^2)\Bigr)+O(\lambda^2)\,.
\]
Therefore, at least to the leading order, we can handle the problem by defining
the renormalized field $[\varphi^2]_R$,
\begin{equation}\label{eq:phi2R}
\varphi^2=Z_{\varphi^2}^{-1}\,[\varphi^2]_R\,,
\end{equation}
where
\[
Z_{\varphi^2}^{-1}=1-\frac\lambda2\,I(0)+O(\lambda^2)\,,\qquad
I(0)=\frac{\Gamma(2-d/2)}{(4\pi)^{d/2}}\,(m^2)^{d/2-2}\,.
\]
Then
\[
\Gamma^{(\varphi^2,2)}=Z_{\varphi^2}^{-1}\,\Gamma^{([\varphi^2]_R,2)}\,,
\]
where the vertex function of the renormalized field $[\varphi^2]_R$ is finite,
to the order $\lambda$.

It is convenient to invert the above relation,
\begin{equation}\label{eq:phi2Rinv}
[\varphi^2]_R(x)=Z_{\varphi^2}\,\varphi^2=Z_{\varphi^2}\,Z^{-1}\,\varphi_0^2\,.
\end{equation}
(Note that in the ratio \eqref{eq:Gphi2} defining $\Gamma^{(\varphi^2,2)}(p_1,p_2)$
the field renormalization factor $Z$ cancels out.) As with other renormalization
parameters, we assume that $Z_{\varphi^2}$ expands in power series in $\lambda$,
\[
Z_{\varphi^2}=z^{(0)}_{\varphi^2}+\lambda\,z^{(1)}_{\varphi^2}+\lambda^2\,z^{(2)}_{\varphi^2}+\dots
\]
In the above calculation we have set $z^{(0)}_{\varphi^2}=1$; with this we have
\[
z^{(1)}_{\varphi^2}=\frac12\,I(0)+\text{finite}\,.
\]
By inserting the renormalized field $[\varphi^2]_R(x)$ instead of $\varphi^2(x)$
we replace the $\varphi^2$ vertex with
\[
\insv[$\varphi^2$]{\node[vtx,minimum size=4.5pt] at (0,0) {};}=1\ \longrightarrow\
\insv{\node[circle,draw,fill=white,inner sep=0pt,minimum size=6pt] at (0,0) {}; \fill (0,0) circle (0.7pt);}=Z_{\varphi^2}\,,
\]
In the perturbation theory (i.e.\ expansion in $\lambda$), this new vertex is
represented as the series in $\lambda$,
\[
\insv{\node[circle,draw,fill=white,inner sep=0pt,minimum size=6pt] at (0,0) {}; \fill (0,0) circle (0.7pt);}
=\insv{\node[circle,draw,fill=white,inner sep=0pt,minimum size=6pt] at (0,0) {}; \fill (0,0) circle (0.7pt); \node[font=\tiny] at (0,0.6) {(0)};}
+\insv{\node[circle,draw,fill=white,inner sep=0pt,minimum size=6pt] at (0,0) {}; \fill (0,0) circle (0.7pt); \node[font=\tiny] at (0,0.6) {(1)};}
+\insv{\node[circle,draw,fill=white,inner sep=0pt,minimum size=6pt] at (0,0) {}; \fill (0,0) circle (0.7pt); \node[font=\tiny] at (0,0.6) {(2)};}
+\cdots
\]
where
\[
\insv{\node[circle,draw,fill=white,inner sep=0pt,minimum size=6pt] at (0,0) {}; \fill (0,0) circle (0.7pt); \node[font=\tiny] at (0,0.6) {$(n)$};}
=\lambda^n\,z^{(n)}_{\varphi^2}
\]
In computing the correlation functions with the insertion of $[\varphi^2]_R$ one
has to take into account these new field renormalization vertices along with
the counterterm vertices associated with the counterterm part of the action. The
field renormalization constants $z^{(n)}_{\varphi^2}$ are to be adjusted, order
by order in the perturbation theory, to cancel the divergent parts of the
diagrams.

As in the case of the renormalizations of the parameters of the action, there
is freedom of finite renormalizations
\[
[\varphi^2]_R\to Z_{\text{finite}}\,[\varphi^2]_R\,.
\]
This ambiguity can be fixed by imposing some normalization condition. For
instance, one can fix the normalization of the field $[\varphi^2]_R$ by
\[
-\Gamma^{([\varphi^2]_R,2)}(0,0)=1\,,
\]
where $-\Gamma^{([\varphi^2]_R,2)}(p_1,p_2)$ is the momentum-space amputated
correlation function
\[
-\Gamma^{([\varphi^2]_R,2)}(p_1,p_2)\,W(p_1)\,W(p_2)=\text{f.t.}\,\Bigl\langle\frac12\,[\varphi^2]_R(x)\,\varphi(x_1)\varphi(x_2)\Bigr\rangle_c
\]
Under this condition we have
\[
z^{(0)}_{\varphi^2}=1\,,\qquad z^{(1)}_{\varphi^2}=\frac12\,I(0)\,.
\]
Renormalizations of higher composite fields follow the same idea but are more
complicated. In particular, it generally involves the phenomenon of ``operator
mixing'': renormalizations of higher composite fields may require forming
linear combinations of them, with renormalization factors becoming matrices.
Brief discussion of the simplest case of this phenomenon, mixing of the fields
$\varphi^3$ and $\Delta\varphi$, is given below. More systematically this
phenomenon is discussed in the context of \emph{Renormalization Group}.

\subsection*{Higher composite fields. Operator mixing}

Simple multiplicative renormalization of the field $\varphi^3$, i.e.\ the form
$[\varphi^3]_R=Z_3(\epsilon)\,\varphi^3$, is not sufficient for removing
divergences in all correlation functions. Instead, one has to look for the
renormalized field in the form of the linear combination
\begin{equation}\label{eq:phi3R}
\frac1{3!}\,[\varphi^3]_R=\frac{Z_3}{3!}\,\varphi^3+Z_{3,1}\,\varphi-Z_{3,1''}\,\partial_\mu^2\varphi\,.
\end{equation}
The reason for adding the term $\sim\partial^2\varphi$ will become clear in a
moment. As usual, the coefficients $Z$ are assumed to be certain power series in
the renormalized coupling constant,
\begin{align*}
Z_3&=1+\lambda\,z^{(1)}_3+\cdots\,,\\
Z_{3,1}&=z^{(0)}_{3,1}+\lambda\,z^{(1)}_{3,1}+\cdots\,,\\
Z_{3,1''}&=\lambda\,z^{(1)}_{3,1''}+\cdots\,.
\end{align*}
Therefore, insertion of the renormalized field $\frac1{3!}[\varphi^3]_R$ gives
rise to series of vertices,
\[
\dia[0.4]{\draw[prop,wavy] (-2.2,0) -- (0,0); \draw[prop] (0,0) -- (1.4,1.4); \draw[prop] (0,0) -- (1.8,0); \draw[prop] (0,0) -- (1.4,-1.4);
  \node[vtx,minimum size=4.5pt] at (0,0) {}; \node[font=\tiny] at (-2.6,0.45) {$\varphi^3$};}=1\,,\qquad
\dia[0.4]{\draw[prop,wavy] (-2.2,0) -- (0,0); \draw[prop] (0,0) -- (1.4,1.4); \draw[prop] (0,0) -- (1.8,0); \draw[prop] (0,0) -- (1.4,-1.4);
  \node[otimesv,minimum size=5pt] at (0,0) {}; \node[font=\tiny] at (0,0.6) {(1)}; \node[font=\tiny] at (-2.6,0.45) {$\varphi^3$};}=\lambda\,z^{(1)}_3\,,\qquad\text{etc}
\]
and
\[
\dia[0.4]{\draw[prop,wavy] (-2.2,0) -- (0,0); \draw[prop] (0,0) -- (2,0); \node[otimesv,minimum size=5pt] at (0,0) {}; \node[font=\tiny] at (0,0.6) {(0)}; \node[font=\tiny] at (-2.6,0.45) {$\varphi^3$};}=z^{(0)}_{3,1}\,,\qquad
\dia[0.4]{\draw[prop,wavy] (-2.2,0) -- (0,0); \draw[prop] (0,0) -- (2,0); \node[otimesv,minimum size=5pt] at (0,0) {}; \node[font=\tiny] at (0,0.6) {(1)}; \node[font=\tiny] at (-2.6,0.45) {$\varphi^3$};}
=\lambda\,\bigl(z^{(1)}_{3,1}+p^2\,z^{(1)}_{3,1''}\bigr)\,,\qquad\text{etc.}
\]
Then, for the vertex $\Gamma^{(\varphi^3_R,1)}$ we have the following series of
diagrams
\begin{multline*}
\underbrace{\Bigl(\dia[0.4]{\draw[prop,wavy] (-2.2,0) -- (0,0); \draw[prop,densely dashed] (0,0) -- (1.8,0); \draw[prop] (0,0.6) circle (0.6);
  \node[vtx,minimum size=4.5pt] at (0,0) {}; \node[font=\tiny] at (-2.6,0.45) {$\varphi^3$};}
+\dia[0.4]{\draw[prop,wavy] (-2.2,0) -- (0,0); \draw[prop,densely dashed] (0,0) -- (1.8,0); \node[otimesv,minimum size=5pt] at (0,0) {}; \node[font=\tiny] at (0,0.6) {(0)}; \node[font=\tiny] at (-2.6,0.45) {$\varphi^3$};}\Bigr)}_{\sim\lambda^0}\\
+\underbrace{\Bigl(\dia[0.4]{\draw[prop,wavy] (-2.2,0) -- (0,0); \draw[prop] (0,0) -- (2,0); \draw[prop] (1,0) ellipse (1 and 0.55); \draw[prop,densely dashed] (2,0) -- (3.2,0);
  \node[vtx,minimum size=4.5pt] at (0,0) {}; \node[vtx] at (2,0) {}; \node[font=\tiny] at (-2.6,0.45) {$\varphi^3$};}
+\dia[0.4]{\draw[prop,wavy] (-2.2,0) -- (0,0); \draw[prop,densely dashed] (0,0) -- (1.8,0); \draw[prop] (0,0.6) circle (0.6);
  \node[otimesv,minimum size=5pt] at (0,0) {}; \node[font=\tiny] at (0,0.6) {(1)}; \node[font=\tiny] at (-2.6,0.45) {$\varphi^3$};}
+\dia[0.4]{\draw[prop,wavy] (-2.2,0) -- (0,0); \draw[prop,densely dashed] (0,0) -- (1.8,0); \node[otimesv,minimum size=5pt] at (0,0) {}; \node[font=\tiny] at (0,0.6) {(1)}; \node[font=\tiny] at (-2.6,0.45) {$\varphi^3$};}\Bigr)}_{\sim\lambda^1}+\cdots
\end{multline*}
We already know from the free theory that
\begin{equation}\label{eq:z031}
z^{(0)}_{3,1}=-\frac12\,D(0)+\text{finite}=-\frac12\int\frac{d^dk/(2\pi)^d}{k^2+m^2}+\text{finite}\,.
\end{equation}
The finite parts here, and in higher order counterterms are arbitrary in that
they are not constrained by general requirement of locality. After all, finite
renormalization of $[\varphi^3]_R$, or adding finite portions of (already
renormalized) fields $\varphi$ and $\partial^2\varphi$ produces equally good
local field that one can take for renormalized $\varphi^3$. One can fix this
ambiguity by adopting minimal subtraction scheme, or by imposing some
normalization conditions, say
\begin{equation}\label{eq:phi3norm}
-\Gamma^{\varphi^3_R,1}(p^2)=O(p^4)\,,
\end{equation}
in addition to the one we have taken already,
\[
-\Gamma^{\varphi^3_R,3}(0,0,0)=1\,.
\]
With this normalization condition, the finite term in \eqref{eq:z031} must be
zero, so that the two diagrams $\sim\lambda^0$ cancel exactly. At the order
$\lambda^1$, the first of the above diagrams is
\[
\dia[0.4]{\draw[prop,wavy] (-2.2,0) -- (0,0); \draw[prop] (0,0) -- (2,0); \draw[prop] (1,0) ellipse (1 and 0.55); \draw[prop,densely dashed] (2,0) -- (3.2,0);
  \node[vtx,minimum size=4.5pt] at (0,0) {}; \node[vtx] at (2,0) {}; \node[font=\tiny] at (-2.6,0.45) {$\varphi^3$};}
=-\frac\lambda{3!}\,\Pi(p^2)\,,
\]
where $\Pi(p^2)$ is already familiar integral
\[
\Pi(p^2)=\int\frac{d^dk_1\,d^dk_2}{(2\pi)^{2d}}\,\frac1{(k_1^2+m^2)(k_2^2+m^2)\bigl((p+k_1+k_2)^2+m^2\bigr)}\,.
\]
We have found that
\[
\Pi_r(p^2)=\Pi(p^2)-\Pi(0)-p^2\,\Pi'(0)
\]
is finite, and the above normalization condition is satisfied by the choice
\[
z^{(1)}_{3,1}=\frac1{3!}\,\Pi(0)-\frac{z^{(1)}_3}2\,D(0)\,,\qquad
z^{(1)}_{3,1''}=\frac1{3!}\,\Pi'(0)\,,
\]
where $z^{(1)}_3=\frac32I(0)$ is the first-order counterterm associated with
$Z_3$. Under this choice all singular parts cancel, and the above three diagrams
$\sim\lambda^1$ sum to finite quantity $-\frac\lambda{3!}\Pi_r(p^2)$.

We have seen few lectures ago that the fields in the correlation functions
satisfy the equation of motion in the ``weak sense'',
\begin{equation}\label{eq:weakEOM}
\bigl(m_0^2-\partial_\mu^2\bigr)\varphi_0(x)+\frac{\lambda_0}{3!}\,\varphi_0^3(x)\simeq0\,,
\end{equation}
the ``weak'' equation $F(x)\simeq0$ means that
\[
\bigl\langle F(x)\,\mathcal O_1(x_1)\cdots\mathcal O_n(x_n)\bigr\rangle=0\qquad\text{for}\quad x\neq x_1,\dots,x_n\,,
\]
with any local fields $\mathcal O_i$. The above ``weak'' equation was obtained by
formal manipulations with the functional integral, namely by performing the
change of variables
\[
\varphi_0(x)\to\varphi_0(x)+\epsilon(x)
\]
in the functional integral written in its ``bare'' form. It involves therefore
the bare parameters and the bare field $\varphi_0$. As it stands it makes sense
only if some regularization is present. One can repeat the same manipulation with
the functional integral expressed through the renormalized field and
renormalized parameters. This way one obtains
\[
\bigl(m^2+\delta m^2\bigr)\varphi(x)-(1+\delta Z)\,\partial_\mu^2\varphi(x)+\frac{\lambda+\delta\lambda}{3!}\,\varphi^3(x)\simeq0\,,
\]
which of course is just the same ``bare'' equation \eqref{eq:weakEOM} expressed
through renormalized parameters and renormalized $\varphi$. It contains the
singular counterterm coefficients as well as the field $\varphi^3$ which, before
renormalization, is also singular. Again, it makes sense only in regularized
theory. However now we can use \eqref{eq:phi3R} to exclude the field $\varphi^3$
in favor of the renormalized field $[\varphi^3]_R$. One then obtains the
equation
\begin{equation}\label{eq:renEOM}
m^2B_{m^2}\,\varphi(x)-\partial_\mu^2\varphi(x)+\frac\lambda{3!}\,B_\lambda\,[\varphi^3]_R(x)\simeq0\,,
\end{equation}
where
\[
B_{m^2}=\frac1{m^2}\,\frac{(m^2+\delta m^2)\,Z_3-(\lambda+\delta\lambda)\,Z_{3,1}}{(1+\delta Z)\,Z_3-(\lambda+\delta\lambda)\,Z_{3,1''}}\,,\qquad
B_\lambda=\frac{1+\delta\lambda/\lambda}{(1+\delta Z)\,Z_3-(\lambda+\delta\lambda)\,Z_{3,1''}}\,.
\]
The equation \eqref{eq:renEOM} states certain relation between the correlation
functions of fully renormalized fields which remain finite when the
regularization is removed (i.e.\ when $\epsilon$ is set equal 0). This suggests
that the coefficients $B_{m^2}$ and $B_\lambda$, which appear as complicated
combination of singular counterterms, are in fact finite. This finiteness is
direct consequence of the existence of renormalized perturbation theory. Let us
check it in the leading order in $\lambda$. As $\delta\lambda/\lambda=a^{(2)}\lambda+O(\lambda^2)$,
$Z_3=1+z_3^{(1)}\lambda+O(\lambda^2)$ and $\delta Z=O(\lambda^2)$ we find for
example
\[
B_\lambda=1+\bigl(a^{(2)}-z^{(1)}_3\bigr)\,\lambda+O(\lambda^2)\,.
\]
It is not difficult to check that the difference $a^{(2)}-z^{(1)}_3$ is finite
in any renormalization scheme.

Note that canonical mass dimension of $\varphi^3$ at $d=4$ is
\[
[\varphi_0^3]=3\,,\qquad\text{and}\qquad[\varphi]=1\,,\qquad[\partial^2\varphi]=3\,.
\]
The renormalized field $[\varphi^3]_R$ is a linear combination of the bare fields
of canonical dimension 3 or lower, which have the same symmetries as
$\varphi^3$. It is not difficult to check that there are no other composite
fields, besides $\varphi_0$, $\partial^2\varphi_0$, and $\varphi_0^3$ itself,
which i) have canonical dimension $\le3$, ii) transform as scalars under
coordinate rotations, and iii) change sign under obvious symmetry of the
$\varphi^4$ theory: $\varphi\to-\varphi$. This is a general feature of
renormalization of composite fields. In order to eliminate divergences, it is
sufficient to look for renormalized fields in the forms of linear combinations of
the bare fields of the same or lower canonical dimensions, and of the same
symmetries. In doing so one usually assumes that the fields of lower dimensions
are already renormalized, and so renormalizations of composite fields goes step
by step, from lower dimensions up.

To illustrate this further, let us briefly discuss renormalization of the fields
of canonical dimension 4. There are two new fields of this dimension,
\begin{equation}\label{eq:dim4ops}
\mathcal O_I=\frac12\,(\partial\varphi)^2\,,\qquad\mathcal O_{II}=\frac1{4!}\,\varphi^4\,.
\end{equation}
Along with these we must consider two other fields
\[
\mathcal O_{III}=\frac12\,\partial^2\varphi^2\,,\qquad\mathcal O_{IV}=\frac12\,\varphi^2\,,
\]
because the first of them also has canonical dimension 4, and canonical
dimension of the second is 2 which is lower than 4, and both have the same
$\varphi\to-\varphi$ symmetry. Of course renormalization of these last two fields
is completely determined by the renormalization of the field $\varphi^2$ which we
already considered. At the same time renormalized fields corresponding to
\eqref{eq:dim4ops} must be obtained as linear combinations of the bare fields
$\mathcal O_{A,0}$, $A=I,II,III,IV$, i.e.
\[
\mathcal O_{A,R}=\sum_B\bigl[Z^{-1}\bigr]_{AB}\,\mathcal O_{B,0}\,,
\]
where $Z$ is $4\times4$ matrix
\[
Z=\begin{pmatrix}\mathbf A&\mathbf B\\0&\mathbf C\end{pmatrix}\,.
\]
Here $\mathbf C=\operatorname{diag}(Z_{\varphi^2},Z_{\varphi^2})$, and $2\times2$
matrices $\mathbf A$ and $\mathbf B$ must be found from suitable normalization
conditions order by order in $\lambda$. Note that this matrix has triangular (in
terms of dimensions) form: fields of higher dimensions do not mix into the fields
of lower dimensions. Renormalizations of composite fields of yet higher canonical
dimensions, and renormalizations of tensor fields, follow the same pattern.
